\chapter*{Introduction}
\label{ch:introduction}
Cancer is an evolutionary process. A tumour develops from a single founding cell and
grows through repeated cycles of cell division, and at each division somatic
mutations arise. Most of these mutations are neutral, leaving the cell's growth and
fitness unchanged; a few give a selective advantage and seed new subclones that
expand on their own. A tumour observed at the moment of sequencing is therefore a
mixture of clones, each carrying its own history of mutations and each having grown
at its own pace.

Reconstructing that history, meaning which events happened, in what order, and how
quickly each clone expanded, is one of the central problems of cancer data science.
The timing of mutations and copy-number changes, and the speed at which clones grow,
say something about how a tumour has progressed and, possibly, about how it will
behave, so a reconstruction that is both accurate and aware of its own uncertainty
is useful scientifically and clinically. The difficulty is that the process is
observed only once: a single snapshot has to be used to infer something that
unfolded over time. This makes careful statistical modelling, and a clear statement
of what can and cannot be recovered from one observation, essential.

The data for this task come in two main forms, which differ in cost, resolution and
the kind of information they give. Bulk sequencing is still comparatively cheap, but
it returns the average signal from a mixed population of cells and so loses
information about the sizes of the clonal subpopulations, which limits how well the
diversity of the tumour can be resolved. Single-cell sequencing instead exposes
tumour heterogeneity in much greater detail and gives information about the subclonal
populations that bulk experiments cannot see. This resolution comes at a price:
single-cell reads are typically noisy and harder to interpret, so the data need
careful processing.

A first step in analysing such data is to reconstruct the phylogenetic tree relating
the sampled cells. Several methods do this, for instance \textit{Pharming}
\parencite{weber2026pharming} and \textit{BiTSC}$^2$ \parencite{chen2022bitsc2},
but a tree inferred from data is an estimate, not a ground truth:
it fixes the ordering constraints among events without, by itself, telling us when
those events occurred or how fast the clones grew. This is the gap the present work
addresses. Taking the inferred phylogeny as its starting point, the thesis develops
a Bayesian hierarchical model that estimates the timing of mutation and
copy-number events together with the clone-specific growth rates. The model ties
several signals to the same latent timeline: the number of mutations accumulated
along each branch, the way mutations in a copy-number region split into those arising
before and after the alteration, and the relative abundance of the clones. It places
a hierarchical prior on growth, so that closely related clones share statistical
strength while distinct ones remain free to differ.

A comparable amount of attention is given to showing that these inferences can be
trusted. The model is validated in stages of increasing difficulty. It is first
fitted to data drawn from its own generative process, where parameter recovery and
simulation-based calibration check that the computation is faithful and the reported
uncertainties are calibrated; it is then tested on synthetic data from ProCESS, a
spatial tumour-evolution simulator, which shows how the inference behaves in a more realistic setting. This staged validation is itself part of
the contribution of the thesis.

The dissertation is organized as follows. Chapter~1 reviews the biological background
and the sequencing technologies, and frames the inference problem in statistical
terms. Chapter~2 develops the models: a first, simpler model that treats the growth
rates as known and infers only the event times, and a second that estimates the
growth rates as well. Chapter~3 presents the simulation study: the ideal
data-generating process, the calibration analysis, and the experiments on ProCESS
data. Chapter~4 discusses the findings, their limitations, and directions for future
work. Two appendices collect the mathematical derivations underlying the model and
the details of its computational implementation.
